\section{Recomendaciones prácticas: cuándo usar cada tipo de plan}

Combinando el contenido de esta clase, se puede repartir la función de los dos mecanismos así:

\begin{table}[H]
\centering
\begin{tabular}{lll}
\toprule
Mecanismo & Uso principal & Escenario típico \\
\midrule
Plan Mode & Aclarar requisitos, formar la especificación & El objetivo del usuario es ambiguo, hay varias rutas posibles para la tarea, se necesita elegir entre alternativas \\
\texttt{update\_plan} & Gestión del avance durante la ejecución & Ya se comenzó la tarea, se necesita mostrar el estado y ajustarlo dinámicamente \\
todo en archivo & Persistencia entre sesiones & Tareas largas, tareas por lotes, se necesita recuperación y auditoría \\
\bottomrule
\end{tabular}
\caption{Reparto de funciones entre los tres soportes de planificación.}
\end{table}

Para un coding agent, el flujo de trabajo más estable suele ser: primero explorar la base de código, y si es necesario entrar a Plan Mode para aclarar la especificación; durante la ejecución usar \texttt{update\_plan} para gestionar el avance actual; y si la tarea abarca un periodo muy largo, mantener además un plan en archivo.

\begin{warningbox}{Un plan no es mejor entre más detallado}
El valor de un plan está en reducir la incertidumbre, no en generar una sensación de formalidad. Un plan demasiado detallado hace que el agent desperdicie contexto en pasos de bajo valor; un plan demasiado general, en cambio, no logra guiar la ejecución. Un buen plan debe capturar únicamente los pasos que en verdad afectan el orden de ejecución y el control de riesgos.
\end{warningbox}

\subsection{Resumen de la sección}

Plan Mode, \texttt{update\_plan} y el todo en archivo son tres niveles. Cada uno resuelve, respectivamente, el problema de aclarar la colaboración, el del estado de ejecución y el de la persistencia a largo plazo.

